% SEGMENT OLP-0688-S01
% Part: proof-theory
% Chapter: propositions-as-types
% Section: proof-terms

\documentclass[../../../include/open-logic-section]{subfiles}

\begin{document}

\olfileid{int}{pty}{ter}

\olsection{நிறுவல் சொற்கள்}

\Log{N1i} இல், $!A^x$ என்பது போல, அனுமானங்களுக்கு
மாறிகளை விடுவிப்பு அடையாளங்களாக இடுகிறோம்.
\Log{N2i} இல், ஒவ்வொரு $\Gamma \Sequent !A$ தொடருருவின்
இடப்புறச் சூழலான~$\Gamma$ இல் மாறி--!!{formula} இணைகளின்
பட்டியல் உள்ளது. ஆகவே, \Log{N2i} இல் !!a{proof} ஒன்றின்
ஒவ்வொரு தொடருருவும், இதுவரை அந்த !!{proof} என்ன
!!{prove} செய்தது ($!A$) என்பதோடு, அந்நிலையில் எந்த
அனுமானங்கள் திறந்துள்ளன ($\Gamma$) என்பதையும் தருகிறது.
$!A$ \emph{எவ்வாறு} !!{prove} செய்யப்பட்டது என்ற
தகவலையும் ஒவ்வொரு தொடருருவில் சேர்த்தால் என்ன?
இதற்காக, $\Gamma \Sequent N : !A$ என்ற வடிவத் தொடருருக்களைக்
கொண்ட \emph{சொல்-குறியிட்ட} \Log{tN2i} முறைமையை
வரையறுக்கலாம். இங்கு $\Gamma$ உம் $!A$ உம் முன்போலவே:
முந்தையது $x: !B$ வடிவிலான இணைகளின் கணம்;
பிந்தையது அத்தொடருருவில் முடியும் துணை-!!{proof},
$\Gamma$ விலுள்ள அனுமானங்களிலிருந்து பின்வருகிறது
என்று நிறுவிய !!{formula}. $N$ என்பது அத்தொடருருவில்
முடியும் !!{proof} ஐக் குறியாக்கும் ஒரு சொல்லாக இருக்கும்.

% SEGMENT OLP-0688-S02
ஒவ்வொரு தொடருருவுக்கும் நாம் அளிக்கும் சொற்கள்,
$x$, $y$, \dots{} போன்ற மாறிகளிலிருந்து, இதுவரை
செய்யப்பட்ட அனுமானங்களைப் பதிவுசெய்யும் சார்புச்
செயற்குறிகளைக் கொண்டு கட்டப்படுகின்றன. ஒவ்வோர் அனுமான
விதிக்கும் தனிச் சார்புச் செயற்குறி வேண்டும். !!{introduction}
விதிகளுக்குரிய செயற்குறிகளை ``கட்டமைப்பிகள்'' என்றும்
!!{elimination} விதிகளுக்குரியவற்றை ``பகுப்பிகள்'' என்றும்
அழைக்கிறோம். ஒவ்வொரு செயற்குறியும், அதற்குரிய விதியில்
எத்தனை முன்கூற்றுகள் உள்ளனவோ அத்தனை உள்ளீடுகளைப்
பெறும். எடுத்துக்காட்டாக, \Intro{\land}, \Elim{\land}[1]
ஆகியவற்றுக்குரிய செயற்குறிகள் முறையே $c^\land$
(இரு உள்ளீடுகள்), $d^\land_1$ (ஓர் உள்ளீடு) ஆகலாம்.
அப்போது \Log{tN2i} இல் இவ்விதிகள் பின்வருமாறு அமையும்:
\[
\Axiom$\Gamma_1 \fCenter N: !A$
\Axiom$\Gamma_2 \fCenter M: !B$
\RightLabel{\Intro{\land}}
\BinaryInf$\Gamma_1, \Gamma_2 \fCenter c^\land(N, M) : !A \land !B $
\DisplayProof
\quad
\Axiom$\Gamma \fCenter N: !A \land !B$
\RightLabel{\Elim{\land}[1]}
\UnaryInf$\Gamma \fCenter d^\land_1(N) : !A$
\DisplayProof
\]

% SEGMENT OLP-0688-S03
ஓர் அனுமான விதி ஓர் அனுமானத்தை விடுவிக்குமானால்,
அதாவது அதன் ஒரு முன்கூற்றில் $x:!A$ என்று குறிக்கப்பட்ட
!!{formula} இருந்தும் முடிவின் சூழலில் அது இல்லாவிட்டால்,
இதை நிறுவல் சொல்லிலும் காட்ட வேண்டும். அத்தகைய
விதிகளுக்குரிய செயற்குறிகள், தொடர்புடைய மாறிகளைக்
\emph{கட்டுப்படுத்துகின்றன}. எடுத்துக்காட்டாக,
\Intro{\lif} விதி, $x:!A, \Gamma \Sequent !B$ என்ற
முன்கூற்றிலிருந்து $\Gamma \Sequent !A \lif !B$ என்ற
முடிவிற்குச் செல்கிறது. $c^{\lif}$ என்ற கட்டமைப்பி
ஓர் உள்ளீட்டைப் பெறுகிறது; அத்துடன் $x$ கட்டுப்பட்ட
மாறியாகிறது என்பதையும் தெரிவிக்க வேண்டும். $x$ என்ற
அடையாளமுள்ள அனுமானம் விடுவிக்கப்பட்டதை நிறுவல்
சொல் இவ்வாறு வெளிப்படுத்துகிறது. முன்கூற்றின் நிறுவல்
சொல்~$N$ என்றால், முடிவின் நிறுவல் சொல்லை
$c^{\lif}([x]N)$ எனவும் விதியைப் பின்வருமாறும் எழுதலாம்:
\[
\Axiom$x:!A, \Gamma \fCenter N: !B$
\RightLabel{\Intro{\lif}}
\UnaryInf$\Gamma \fCenter c^{\lif}([x]N) : !A \lif !B$
\DisplayProof
\]

% SEGMENT OLP-0688-S04
!!a{proof} ஐ முழுமையாக மீட்டமைக்க, இன்னும் சில
தகவல்கள் தேவை. ஒரு விதியின் முன்கூற்று !!{formula}s
மட்டுமே அதன் முடிவு !!{formula} எது என்பதை முழுவதுமாகத்
தீர்மானிக்கவில்லை என்றால், அந்தத் தகவலைச் சொல்லுடன்
சேர்க்க வேண்டும். \Intro{\lif} இல் இவ்வாறு இருப்பதால்,
$c^{\lif}([x:!A]N)$ என எழுதுவோம். \Intro{\lor} இலும்
\FalseInt{} இலும் அவ்வாறே இருப்பதால், தேவையான தகவலைத்
தாங்கும் செயற்குறிகளைப் பயன்படுத்துவோம். எடுத்துக்காட்டாக,
\[
\Axiom$\Gamma \fCenter N:!A$
\RightLabel{\Intro{\lor}[1]}
\UnaryInf$\Gamma \fCenter c^{\lor{!B}}_1(N):!A \lor !B$
\DisplayProof
\quad
\Axiom$\Gamma \fCenter N:\lfalse$
\RightLabel{\FalseInt}
\UnaryInf$\Gamma \fCenter d^\lfalse_{!A}(N):!A$
\DisplayProof
\]

% SEGMENT OLP-0688-S05
இந்தக் கருத்துகளைப் பயன்படுத்தி, ஒவ்வொரு தொடருருவும்
$\Gamma \Sequent N : !A$ என்ற வடிவில் இருக்கும்
தொடருரு-முறை இயல்பான வருவித்தல் முறைமையைக் கட்டலாம்;
இங்கு $N$ ஒரு \emph{நிறுவல் சொல்}.

இத்தகைய நிறுவல் சொற்களுக்கும் \emph{வகையிட்ட லாம்டா
கலனம்} எனப்படும் கலனத்தின் சொற்களுக்கும் நெருங்கிய
தொடர்பு உள்ளது. அதற்கேற்ப, மேலுள்ள பொதுவான
கட்டமைப்பி, பகுப்பிச் செயற்குறிகளுக்குப் பதிலாக,
லாம்டா கலனம் பற்றிய ஆய்வுகளில் வழக்கிலுள்ள
செயற்குறிகளைப் பயன்படுத்துவோம். எடுத்துக்காட்டாக,
$c^{\lif}([x:!A]N)$ என்பதற்குப் பதில்
$\lambd[\typeof{x}{!A}][N]$ என்றும்,
$d^{\lif}(N,M)$ என்பதற்குப் பதில் $(NM)$ என்றும்,
$c^\land(N,M)$ என்பதற்குப் பதில் $\pair{N}{M}$ என்றும்,
$d^\land_i(N)$ என்பதற்குப் பதில் $\proj{i}{N}$ என்றும்
எழுதுவோம்.

% SEGMENT OLP-0688-S06
\begin{defn}
  நிறுவல் சொற்கள் பின்வரும் விதிகளால் கட்டமைப்புத்
  தொகுத்தறிதலின்படி உருவாக்கப்படுகின்றன:
  \begin{enumerate}
  \item ஒவ்வொரு மாறி~$x$ உம் ஒரு நிறுவல் சொல்
    (அடிக்கோள்கள்).
  \item $x$ ஒரு மாறியாகவும், $!A$ ஓர் !!a{formula} ஆகவும்,
    $N$ ஒரு நிறுவல் சொல்லாகவும் இருந்தால்,
    $\lambd[\typeof{x}{!A}][N]$ உம் ஒரு நிறுவல் சொல்
    (\Intro\lif).
  \item $N$ உம் $M$ உம் நிறுவல் சொற்கள் என்றால்,
    $(NM)$ உம் ஒரு நிறுவல் சொல் (\Elim\lif).
  \item $N$ உம் $M$ உம் நிறுவல் சொற்கள் என்றால்,
    $\pair{N}{M}$ உம் ஒரு நிறுவல் சொல் (\Intro\land).
  \item $N$ ஒரு நிறுவல் சொல் என்றால், $i$ என்பது
    $1$ அல்லது~$2$ ஆக இருக்கையில், $\proj{i}{N}$ உம்
    ஒரு நிறுவல் சொல் (\Elim\land[i]).
  \item $N$ ஒரு நிறுவல் சொல்லாகவும் $!A$ ஒரு வாய்பாடாகவும்
    இருந்தால், $i$ என்பது $1$ அல்லது~$2$ ஆக இருக்கையில்,
    $\inj[!A]{i}{N}$ உம் ஒரு நிறுவல் சொல்
    (\Intro\lor[i]).
  \item $M$, $N_1$, $N_2$ நிறுவல் சொற்களாகவும் $x_1$,
    $x_2$ மாறிகளாகவும் இருந்தால்,
    $\dcase{M}{x_1}{N_1}{x_2}{N_2}$ ஒரு நிறுவல் சொல்
    (\Elim\lor).
  \item $N$ ஒரு நிறுவல் சொல்லாகவும் $!A$ ஓர் !!a{formula}
    ஆகவும் இருந்தால், $\abort{!A}{N}$ ஒரு நிறுவல் சொல்
    (\FalseInt).
  \end{enumerate}
\end{defn}

% SEGMENT OLP-0688-S07
ஒவ்வொரு நிறுவல் சொல்லும் ஒரு !!a{proof} இன்
நிறுவல் சொல்லாக அமையாது. எடுத்துக்காட்டாக,
$\pair{N}{M}$ என்பது \Intro{\land} அனுமான விதியில்
முடியும் ஒரு நிறுவலைக் குறிக்கிறது; ஆகவே அதன் முடிவு
!!{formula} ஒரு இணையலான $!A \land !B$. இது
\Elim{\lif} விதியின் முதன்மை முன்கூற்றாக அமைய
முடியாது. எனவே $(\pair{N}{M}P)$ என்பது சரியான
நிறுவலின் நிறுவல் சொல்லாக முடியாது.
\Log{tN2ip} இன் விதிகளுக்குப் பொருந்தும் நிறுவல்
சொற்கள் மட்டுமே \emph{சரியான} நிறுவல் சொற்கள்.
அவ்விதிகள் \olref{tab:tN2ip} இல் தரப்பட்டுள்ளன.

\subfile{rules-tN2}

\end{document}
